%% P25 --- Sumy zmiennych losowych: centralne twierdzenie graniczne i metoda Monte Carlo
%%
%% Zalążek, wygenerowany przez tools/gen_stubs.py z tools/programs.json.
%% Poniższy opis jest kontraktem tego programu: co ma uzasadnić i co ma dać
%% czytelnikowi. Napisanie programu oznacza usunięcie bloku \programstub{} i
%% zastąpienie go programem -- po czym ten plik nie jest już generowany.

\program{Sumy zmiennych losowych: centralne twierdzenie graniczne i metoda Monte Carlo}
\label{prog:P25}

\programstub{%
\textit{[Opis do przetłumaczenia --- patrz wydanie angielskie.]}

Argument: variances of independent quantities add; averages concentrate; the
error of an average falls as \code{1/sqrt(n)} and that single rate governs
an astonishing amount of practice. Payoff: **the derivation of the
$1/sqrt(d_k)$ in attention.** A dot product of two $d_k$-dimensional vectors
with independent unit-variance entries has variance $d_k$, so its standard
deviation grows as $sqrt(d_k)$; feed that into a softmax and it saturates,
and the gradient dies; divide by $sqrt(d_k)$ and the logits have unit
variance regardless of head size. The scaling is a variance correction and
nothing else, and the reader derives it here and measures it in P31. Also:
Monte Carlo error is the same \code{1/sqrt(n)}, so the number of samples
needed to halve an error bar is four times as many \dash{} the fact that
prices every evaluation run in the book.

ADDITION, from the review: nothing in the book explained why a training run
diverges at initialisation. Extend this program with signal propagation --
the variance through a linear layer as fan-in times weight variance, the
scale that holds it at one, and what a non-linearity does to it, which is
where the factor of two in He initialisation comes from. It belongs here
because it is a variance argument, and it is the single most common cause of
a run that never starts.

\medskip
\textbf{Planowana długość:} 55 ramek.
}
